\chapter[Information Functionals]{Information Functionals on the WIS Manifold}\label{ch:sec:information-functionals-on-the-wis-manifold}

This chapter provides the unified classification of security functionals---smooth, piecewise-smooth, and rank-based---that organises the remainder of Part~IV. The Lipschitz bounds of Chapters~15--16 apply uniformly across this classification, and the piecewise-smooth theory of Chapters~17--18 extends them to non-differentiable measures.


\section{Introduction}\label{ch:sec:introduction-functionals}
We have established that the UOD provides a canonical geometric distance
on the posterior manifold.  But cryptography and information theory use
many different measures---Shannon entropy, R\'enyi entropy, min-entropy,
Bayes vulnerability, guessing entropy.  Do we need a separate geometric
theory for each?  Or do they all fit into a single framework?  This
chapter shows that every information or security functional---whether
smooth or piecewise-smooth---is a scalar field on the same WIS manifold,
and the UOD bounds their variations universally.

The purpose of this chapter is to develop a unified geometric framework
for information-theoretic and security functionals on the Weighted
Incidence Structure (WIS) manifold introduced in Part~\ref{ch:part:posterior}.
Rather than studying Shannon entropy, R\'enyi entropy, Kullback--Leibler
divergence, or security measures independently, we regard them as scalar
functionals defined on the same posterior manifold
$\mathcal P.$
The WIS construction provides a global coordinate system
$F:\mathcal P\rightarrow\mathbb R^n,$
together with the pullback Euclidean metric
$g=F^*g_E.$
This common geometric structure allows universal analytical results that
are independent of the particular functional under consideration.
\paragraph{Information Functionals.}
An \textbf{information functional} is any mapping
$S:\mathcal P\rightarrow\mathbb R.$
Throughout this chapter we distinguish two broad classes.
\subsection{Smooth information functionals}\label{ch:sec:smooth-information-functionals}
A functional is called \textit{smooth} if
$S\in C^1(\mathcal P),$
or \(C^{2}\) whenever second-order results are required.
Examples include:
\begin{itemize}
\item Shannon entropy
\item Cross entropy
\item Kullback--Leibler divergence
\item R\'enyi entropy
\item Tsallis entropy
\item Collision entropy
\item Jensen--Shannon divergence
\item Differentiable (f)-divergences
\item Mutual information (when represented as a smooth functional on the
    posterior manifold)
\end{itemize}
These functionals admit gradients and Hessians with respect to the WIS
metric.

\subsection{Rank-based information functionals}\label{ch:sec:rank-based-information-functionals}
Rank-based functionals depend on the ordering of posterior probabilities rather than on their exact values.  Typical examples include min-entropy $H_\infty(P)=-\log\max_i P_i$, guessing entropy, Bayes vulnerability, and success probability.  These functionals are piecewise-smooth (see Section~\ref{ch:sec:piecewise-smooth-security-functionals}): they are differentiable on regions where the ranking is fixed, but their gradient changes discontinuously at boundaries where two probabilities cross.
\section{Piecewise-smooth security functionals}\label{ch:sec:piecewise-smooth-security-functionals}
Many security measures are not globally differentiable because they
depend on maxima or rankings.
Typical examples are:
\begin{itemize}
\item Min-Entropy
\item Bayes Vulnerability
\item Guessing Entropy
\item Success Probability
\item Top-(k) Success Probability
\item Rank-based leakage measures
\end{itemize}
Although not globally smooth, these functionals are smooth on regions
where the ordering of probabilities remains unchanged.
\paragraph{WIS Coordinates.}



Part~\ref{ch:part:posterior} established that every posterior distribution admits global
coordinates
$x=F(P).$
Consequently every information functional can be rewritten as
$\widetilde S(x)=S(F^{-1}(x)).$
All subsequent analysis is therefore performed in Euclidean coordinates
while remaining intrinsically defined on the posterior manifold.
\paragraph{Universal Optimality Defect.}
For two posterior distributions (P,Q), define the two-point Universal
Optimality Defect
$\mathcal D(P,Q) = \|F(P)-F(Q)\|^2.$
When one endpoint corresponds to a distinguished optimal posterior
\(P^{\star}\),
$\mathcal D(P) = \mathcal D(P,P^\star).$
The Universal Optimality Defect provides the canonical quadratic
geometric quantity used throughout the remainder of this chapter.



\paragraph{Two Classes of Universal Results.}
The geometric theory naturally separates into two families.



\paragraph{Smooth theory.}



For differentiable functionals:
\begin{itemize}
\item gradient bounds;
\item Hessian bounds;
\item Bregman representations;
\item curvature estimates;
\item quadratic approximations.
\end{itemize}
\paragraph{Piecewise-smooth theory.}



For ranking-based functionals:
\begin{itemize}
\item regularity intervals;
\item ranking transitions;
\item piecewise differentiability;
\item universal geometric bounds on each smooth region.
\end{itemize}
This separation allows a single geometric framework to accommodate
essentially all classical information-theoretic and security measures.
\subsection{Roadmap}\label{ch:sec:roadmap}
Part~\ref{ch:part:functionals} develops this theory as follows.
\begin{itemize}
\item \textbf{Chapter~\ref{ch:sec:universal-geometric-bounds}} establishes universal geometric bounds for smooth
    functionals.
\item \textbf{Chapter~\ref{ch:sec:the-average-hessian-operator-and-universal-bregman-theo}} develops the general Bregman framework and curvature
    theory.
\item \textbf{Chapter~\ref{ch:sec:classical-smooth-information-functionals}} applies the universal results to classical smooth
    information measures.
\item \textbf{Chapter~\ref{ch:sec:piecewise-smooth}} extends the theory to piecewise-smooth security
    functionals.
\item \textbf{Chapter~\ref{ch:sec:universal-security-bounds}} derives unified security bounds from the
    piecewise-smooth structure.
\end{itemize}
\subsection{Perspective}\label{ch:sec:perspective}
The classical approach to information geometry starts from a divergence
and derives a metric.
The WIS construction reverses this philosophy.
A single global geometric structure is established first, after which
information and security functionals appear as scalar fields defined on
the same manifold.
This viewpoint transforms apparently unrelated quantities into different
manifestations of a common geometric framework, preparing the universal
results developed in the subsequent chapters.

\paragraph{Running examples.}
The additional examples illustrate the two regularity classes.  In differential privacy (Section~\ref{ch:sec:differential-privacy-and-database-queries}), the Laplace mechanism produces a smooth functional: the privacy loss $\mathcal L(P\|Q)=\|F(P)-F(Q)\|^2$ is $C^\infty$ on the interior of the simplex, and the Lipschitz bounds of Chapter~\ref{ch:sec:universal-geometric-bounds} apply with constant $L\le 1/\min_i\pi_i$.  Browser fingerprinting (Section~\ref{ch:sec:browser-fingerprinting-and-anonymity}) falls into the piecewise-smooth class: $H_\infty(P)$ depends on the most probable browser and changes discontinuously when the ranking shifts.  Our main running example, the AES S-box, is trivially smooth because $\mathcal D\equiv0$ makes every functional constant.

